\section{The \texttt{treepo} Software}\label{app:treepo-software}

The experiments in Section~\ref{sec:manifesto-deepdive} and the
controlled validations run on one package, \texttt{treepo}. Its design
goal is the operator interface of Section~\ref{sec:framework} made
executable: the framework's objects ($f$, $g$, trees, local-law
probes, audits) are fixed by the package, and everything
setting-specific lives behind a small family contract. The same
harness therefore runs an LLM prompt program, a Fourier neural
operator, and a classical sketch without changing the surrounding
loop, which is what makes the substrate comparisons in the main text
comparisons rather than reconciliations.

\subsection{The Family Contract and \texttt{treepo.fit}}

The public entry point is \texttt{treepo.fit}, which takes a
declarative spec (family, data, schedule) and runs the alternating
$f/g$ ladder of Section~\ref{sec:rile-fg-ladder} for any
\emph{family}: an object implementing the four-method contract
\begin{lstlisting}[language=Python]
class FamilyRuntime(Protocol):
    name: str
    def train_f(self, *, f_init, g, traces, output_dir, iteration): ...
    def train_g(self, *, g_init, f, traces, output_dir, iteration): ...
    def score_roots_with_f(self, *, f, g, trees): ...
    def validate_artifact(self, *, kind, artifact): ...
\end{lstlisting}
A family supplies its own parameterization of the summarizer $g$ and
readout $f$; the surrounding loop (alternation schedule, node
sampling, metrics, audit logging) is family-agnostic. The two families
used in Section~\ref{sec:qsentence-results} implement the contract at
opposite ends of the parameterization spectrum: \texttt{dspy} trains
prompt programs over a fixed open-weight checkpoint by prompt search
against node-level rewards, with text as the inter-node state;
\texttt{fno} embeds leaves once with a small embedding model and
trains Fourier neural operators by gradient descent on node-level
losses, with a learned vector as the state.

\subsection{A Minimal Vignette}

A minimal run is three calls: build the trees, declare the backend,
fit.
\begin{lstlisting}[language=Python]
from treepo import fit
from treepo.tasks.manifesto import (
    make_manifesto_replication_trees, manifesto_prompt_template)

train = make_manifesto_replication_trees(
    split="train", leaf_unit_count=8, doc_unit_kind="qsentence")
evaltrees = make_manifesto_replication_trees(
    split="test", leaf_unit_count=8, doc_unit_kind="qsentence")

result = fit({
    "family": "dspy",              # or "fno"
    "train_data": train,
    "eval_data": evaltrees,
    "backend_config": {
        "model": "openai/nvidia/Gemma-4-31B-IT-NVFP4",
        "prompt_template": manifesto_prompt_template(),
    },
    "axis": {"axis_kind": "leaf_unit_count", "axis_value": 8,
             "max_iterations": 2},
})
print(result.status, result.metrics["internal_f_mae"])
\end{lstlisting}
A complete end-to-end walkthrough (tree construction from raw
platforms, sampling artifacts, fitting, preference-pair export) ships
with the package as
\texttt{examples/methods/run\_manifesto\_end\_to\_end.py}.

\subsection{Research Sweeps}

The research sweeps behind the quasi-sentence tables drive the same
machinery through a CLI that exposes the leaf-size axis and the family
switch; the two substrate legs of the ladder differ in a single flag:
\begin{lstlisting}[language=bash]
# LLM prompt-program leg
python scripts/run_manifesto_qsentence_dspy_ladder.py \
  --family dspy --leaf-qsentences "1,2,4,8,16" \
  --max-iterations 2 --target-dimensions all \
  --dspy-optimizer gepa \
  --dspy-model "openai/nvidia/Gemma-4-31B-IT-NVFP4"

# neural-operator leg: same bundles, same splits
python scripts/run_manifesto_qsentence_dspy_ladder.py \
  --family fno --leaf-qsentences "1,2,4,8,16" \
  --embedding-model google/embeddinggemma-300m \
  --fno-target-dimension rile --fno-epochs 8
\end{lstlisting}

\subsection{Artifact Schema}

Every run leaves behind the same artifact layout regardless of family:
per-stage grid summaries (one row per leaf-size--stage cell, with
internal and external metrics), per-node prediction records (one JSON
line per tree node per evaluation pass, carrying the node's predicted
and gold states), and iteration histories (the full stage trajectory
with the metrics reported in the ladder tables). Cross-family
comparisons such as Table~\ref{tab:qsentence-perdim} are computed by
joining these artifacts on tree and node identifiers; no
family-specific parsing is involved. The saved per-node records are
also what the audit of Section~\ref{sec:estimation} consumes: audit
sampling operates on logged nodes and propensities, so a completed
training run is auditable after the fact without re-running the model.

\subsection{The \texttt{treepo-bench} Classical-Sketch Suite}

The classical validations run through the same \texttt{fit()}
interface via the companion \texttt{treepo-bench classical-sketches}
suite. Beyond the HyperLogLog parity experiment of
Appendix~\ref{sec:hll-parity}, the suite repeats the tree-reduction
comparison for official Apache DataSketches distinct-counting,
frequency, quantile, and set-operation sketches plus Tuple/VarOpt
sampling checks, reporting memory/error curves over capacities and
leaf counts, merge-schedule spread, bound coverage, and distance to
the best official sketch floor. Because the oracle $\fstar$ is a
constructor argument, the same pipeline swaps between an analytic
target (exact cardinality) and an official sketch query function
(the HLL scoring head), which is how the parity and learned rows of
Appendix~\ref{app:classical-parity} are generated from one code path.

\subsection{Reproduction}

Paper assets regenerate via \texttt{paper/ctreepo/regen\_assets.sh};
build instructions for the paper and the Lean formalization are in
\texttt{paper/ctreepo/REPRODUCIBILITY.md}. Each results table in
Appendices~\ref{app:benoit-replication}--\ref{app:gold-seeding} names
the run directory that produced it, and each run directory contains
the exact command line in its job log.
